% ch6_c (书页 263-273 / p278-p288)
%--B-TAIL--
%JOIN%静态的。这是有道理的：黑洞在转动，所以它不是静态的；但它在所有时刻都以完全相同的方式转动，所以它是稳态的。或者也可以这样理解：这个度规不可能是静态的，因为它不具有时间反演不变性——时间反演会把黑洞的角动量反过来。

%FIG{6.6}: {Kerr 度规中使用的椭球坐标 $(r,\theta)$。$r=0$ 是一个二维圆盘；$r=0$ 与 $\theta=\pi/2$ 的交集是这个圆盘边界上的圆环。}

% ---- 书页 263 / p278 ----（本页第一段已在上面 tail 中接续，正文自第二段起）

Kerr 度规还拥有一个 Killing 张量（Killing tensor）。第 3 章中曾把它定义为任何满足如下条件的对称 $(0,n)$ 张量 $\sigma_{\mu_{1}\cdots\mu_{n}}$：
\begin{equation*}
\nabla_{(\lambda}\sigma_{\mu_{1}\cdots\mu_{n})} = 0.
\tag{6.77}
\end{equation*}
在 Kerr 几何中，我们可以定义 $(0,2)$ 张量
\begin{equation*}
\sigma_{\mu\nu} = 2\rho^{2}l_{(\mu}n_{\nu)} + r^{2}g_{\mu\nu}.
\tag{6.78}
\end{equation*}
在这个表达式中，两个矢量 $l$ 和 $n$（升指标后）为
\begin{gather*}
l^{\mu} = \frac{1}{\Delta}\left(r^{2}+a^{2},\, \Delta,\, 0,\, a\right)\\
n^{\mu} = \frac{1}{2\rho^{2}}\left(r^{2}+a^{2},\, -\Delta,\, 0,\, a\right).
\tag{6.79}
\end{gather*}
这两个矢量都是类光的，并且满足
\begin{equation*}
l^{\mu}l_{\mu} = 0,\qquad n^{\mu}n_{\mu} = 0,\qquad l^{\mu}n_{\mu} = -1.
\tag{6.80}
\end{equation*}
有了这些定义，你可以自己验证 $\sigma_{\mu\nu}$ 是一个 Killing 张量。

我们为 Kerr 选定的坐标，使得事件视界出现在那些让 $g^{rr}=0$ 的固定 $r$ 值处。由于 $g^{rr}=\Delta/\rho^{2}$，而 $\rho^{2}\geq 0$，这发生在
\begin{equation*}
\Delta(r) = r^{2} - 2GMr + a^{2} = 0.
\tag{6.81}
\end{equation*}
与 Reissner–Nordström 解一样，这里有三种可能性：$GM>a$、$GM=a$ 和 $GM<a$。最后一种情形的特征是存在裸奇点，而极端情形 $GM=a$ 是不稳定的，正如 Reissner–Nordström 情形一样。由于这些情形在物理上的重要性较低，我们将集中关注 $GM>a$。这时 $\Delta$ 在两个半径处为零，由
\begin{equation*}
r_{\pm} = GM \pm \sqrt{G^{2}M^{2}-a^{2}}.
\tag{6.82}
\end{equation*}
给出。这两个半径都是类光曲面，我们将会看到它们正是事件视界；Kerr 黑洞的侧视图见图 6.7。对这些曲面的分析与 Reissner–Nordström 情形极为类似；不难找到穿过视界延拓的坐标。

由于 Kerr 是稳态而非静态的，$r_{\pm}$ 处的事件视界并不是渐近时间平移 Killing 矢量 $K=\partial_t$ 的 Killing 视界。$K^{\mu}$ 的模由
\begin{equation*}
K^{\mu}K_{\mu} = -\frac{1}{\rho^{2}}\left(\Delta - a^{2}\sin^{2}\theta\right).
\tag{6.83}
\end{equation*}
给出。

% ---- 书页 264 / p279 ----

%FIG{6.7}: {Kerr 解周围的视界结构（侧视图）。事件视界是类光曲面，它标示出一条分界：越过它之后就不可能再返回空间的某个区域。与此不同，静止极限面除了与事件视界相切之处（两极）以外都是类时的；它代表越过之后便无法保持静止观者身份的地方。静止极限面与外侧事件视界之间的能层（ergosphere）是一个可以进入、也可以再离开、但不能在其中保持静止的区域。}

在外侧事件视界处它并不为零；事实上，在 $r=r_{+}$（此处 $\Delta=0$）我们有
\begin{equation*}
K^{\mu}K_{\mu} = \frac{a^{2}}{\rho^{2}}\sin^{2}\theta \geq 0.
\tag{6.84}
\end{equation*}
因此，这个 Killing 矢量在外侧视界处已经是类空的了，只有在北极和南极（$\theta=0,\pi$）处它是类光的。$K^{\mu}K_{\mu}=0$ 的点的轨迹当然就是静止极限面，由
\begin{equation*}
(r-GM)^{2} = G^{2}M^{2} - a^{2}\cos^{2}\theta,
\tag{6.85}
\end{equation*}
给出，而外侧事件视界则由
\begin{equation*}
(r_{+}-GM)^{2} = G^{2}M^{2} - a^{2}.
\tag{6.86}
\end{equation*}
给出。于是，这两个曲面之间存在一个区域，称为\textbf{能层}（ergosphere）。在能层内部，你必须沿着黑洞转动的方向运动（即 $\phi$ 方向）；不过，你仍然可以朝事件视界运动，也可以远离它（要退出能层也没有任何麻烦）。显然，甚至在穿越视界之前，能层就是一个可能发生有趣事情的地方；详情后述。

在急着画共形图之前，我们需要弄清楚真正的曲率奇点的性质；在这个时空中，它并不位于 $r=0$ 处，而是位于 $\rho=0$ 处（曲率不变量 $R_{\rho\sigma\mu\nu}R^{\rho\sigma\mu\nu}$ 在那里发散）。由于 $\rho^{2}=r^{2}+a^{2}\cos^{2}\theta$ 是两个明显非负的量之和，只有当两者同时为零时它才可能为零，即
\begin{equation*}
r = 0,\qquad \theta = \frac{\pi}{2}.
\tag{6.87}
\end{equation*}

% ---- 书页 265 / p280 ----

这个结果看上去有点滑稽，但请记住 $r=0$ 不是空间中的一个点，而是一个圆盘；满足 $r=0$、$\theta=\pi/2$ 的点集实际上就是这个圆盘边缘上的\emph{环}。转动把 Schwarzschild 奇点“软化”了，将它摊开成一个环。

如果你走进环内会发生什么？仔细的解析延拓（我们不打算真的去做）将会揭示，你会到达另一个渐近平直的时空，但它并不是你来自的那个时空的全同副本。新时空由 $r<0$ 的 Kerr 度规描述。结果是 $\Delta$ 永不为零，从而不存在视界。共形图（图 6.8）与 Reissner–Nordström 的情形非常相像，只是现在你可以穿过奇点。由于 Kerr 度规不是球对称的，共形图也就不像前面几种情形那样忠实；图上的一个点代表固定的 $t$ 和 $r$ 值，而对于不同的 $\theta$ 值，它将有不同的几何。

%FIG{6.8}: {$G^{2}M^{2}>a^{2}$ 的 Kerr 解的共形图。与类似的带电解一样，黑洞外部的区域存在着无穷多个副本。}

% ---- 书页 266 / p281 ----

我们不仅有这些经由黑洞与我们所在区域相连的、彼此不同的渐近平直区域所带来的通常怪异性，而且环奇点附近的区域还有额外的病态之处：闭合类时曲线（closed timelike curves，CTC）。如果你考虑一些在 $\phi$ 方向绕行、同时保持 $\theta$ 和 $t$ 固定、$r$ 取很小负值的轨迹，沿这样一条路径的线元是
\begin{equation*}
ds^{2} \approx a^{2}\left(1+\frac{2GM}{r}\right)\mathrm{d}\phi^{2},
\tag{6.88}
\end{equation*}
对于很小的负 $r$，它是负的。由于这些路径是闭合的，它们显然就是 CTC。因此，你有可能在过去遇见你自己，以及随之而来的一切。

当然，我们关于 Kerr 解析延拓所说的一切，都带有我们在讨论 Schwarzschild 和 Reissner–Nordström 时提到过的同样的告诫；真实的引力坍缩不太可能形成这些古怪的时空。尽管如此，拥有精确解总是有用的。此外，对于 Kerr 度规，即使我们停留在事件视界之外，奇怪的事情也在发生，现在我们就转向这些事情。

我们首先更仔细地考察黑洞的角速度。显然，常规的角速度定义必须稍加修改，才能应用于像时空度规这样抽象的对象。让我们考虑在 Kerr 黑洞的赤道面（$\theta=\pi/2$）内、半径 $r$ 处沿 $\phi$ 方向发射的一个光子的命运。它发射的那一瞬间，其动量在 $r$ 和 $\theta$ 方向都没有分量，因此轨迹为类光的条件是
\begin{equation*}
ds^{2} = 0 = g_{tt}\mathrm{d}t^{2} + g_{t\phi}\left(\mathrm{d}t\,\mathrm{d}\phi + \mathrm{d}\phi\,\mathrm{d}t\right) + g_{\phi\phi}\mathrm{d}\phi^{2}.
\tag{6.89}
\end{equation*}
这可以立即解出
\begin{equation*}
\frac{\mathrm{d}\phi}{\mathrm{d}t} = -\frac{g_{t\phi}}{g_{\phi\phi}} \pm \sqrt{\left(\frac{g_{t\phi}}{g_{\phi\phi}}\right)^{2} - \frac{g_{tt}}{g_{\phi\phi}}}.
\tag{6.90}
\end{equation*}
如果在 Kerr 度规的静止极限面上计算这个量，就有 $g_{tt}=0$，两个解为
\begin{equation*}
\frac{\mathrm{d}\phi}{\mathrm{d}t} = 0,\qquad \frac{\mathrm{d}\phi}{\mathrm{d}t} = \frac{a}{2G^{2}M^{2}+a^{2}}.
\tag{6.91}
\end{equation*}
非零解与 $a$ 同号；我们把它解释为光子绕黑洞运动的方向与黑洞转动的方向相同。零解则意味着，沿逆着黑洞转动方向射出的光子在这个坐标系中完全不动。注意，我们并没有给出光子轨迹的完整解，只是证明了它的瞬时速度为零。这就是著名的“参考系拖曳”（dragging of inertial frames）现象的一个例子；
% ---- 书页 267 / p282 ----
它将在第 7 章的一道习题中得到更多探讨。有质量的粒子必定比光子运动得慢，一旦进入静止极限面之内，它们就不可避免地被拖着随黑洞一起转动。当我们向位于 $r_{+}$ 处的外侧事件视界靠近时，这种拖曳仍在继续；我们可以定义事件视界本身的角速度 $\Omega_{\mathrm{H}}$，即粒子在视界处的最小角速度。直接由式 (6.90) 得到
\begin{equation*}
\Omega_{\mathrm{H}} = \left(\frac{\mathrm{d}\phi}{\mathrm{d}t}\right)_{-}(r_{+}) = \frac{a}{r_{+}^{2}+a^{2}}.
\tag{6.92}
\end{equation*}

\section{Penrose 过程与黑洞热力学}

黑洞热力学是广义相对论中最引人入胜、也最神秘的课题之一。不过，为了走到那里，让我们先从一件看上去非常直截了当的事情开始：Kerr 度规中沿测地线的运动。我们知道，如果考虑与 Killing 矢量 $K=\partial_t$ 和 $R=\partial_\phi$ 相联系的守恒量，这样的讨论将会得到简化。就眼前的目的而言，我们可以把注意力限制在有质量的粒子上，对它们可以使用四维动量
\begin{equation*}
p^{\mu} = m\frac{\mathrm{d}x^{\mu}}{\mathrm{d}\tau},
\tag{6.93}
\end{equation*}
其中 $m$ 是粒子的静质量。于是，我们可以取粒子真实的能量和角动量作为我们的两个守恒量，
\begin{equation*}
E = -K_{\mu}p^{\mu} = m\left(1-\frac{2GMr}{\rho^{2}}\right)\frac{\mathrm{d}t}{\mathrm{d}\tau} + \frac{2mGMar}{\rho^{2}}\sin^{2}\theta\,\frac{\mathrm{d}\phi}{\mathrm{d}\tau}
\tag{6.94}
\end{equation*}
以及
\begin{equation*}
L = R_{\mu}p^{\mu} = -\frac{2mGMar}{\rho^{2}}\sin^{2}\theta\,\frac{\mathrm{d}t}{\mathrm{d}\tau} + \frac{m(r^{2}+a^{2})^{2}-m\Delta a^{2}\sin^{2}\theta}{\rho^{2}}\sin^{2}\theta\,\frac{\mathrm{d}\phi}{\mathrm{d}\tau}.
\tag{6.95}
\end{equation*}
这些定义与上一章中使用的守恒量定义有所不同，在上一章里 $E$ 和 $L$ 取的是\emph{每单位质量}的能量和角动量。当然，无论哪种取法，它们都是守恒的。

$E$ 的定义中有一个负号，这是因为在无穷远处 $K^{\mu}$ 和 $p^{\mu}$ 都是类时的，它们的内积为负，而我们希望能量是正的。然而在能层内部，$K^{\mu}$ 变成了类空的；因此我们可以设想有这样的粒子，使得
\begin{equation*}
E = -K_{\mu}p^{\mu} < 0.
\tag{6.96}
\end{equation*}

% ---- 书页 268 / p283 ----

这件事令我们感到的不安，多少可以因如下认识而缓解：只要在静止极限面之外，\emph{所有}粒子的能量都必定为正；因此，能层内部具有负能量的粒子，要么留在能层之中，要么若想逃出去，就必须被加速到能量变为正。

尽管如此，这一认识仍然开辟了一条从转动黑洞中汲取能量的途径；这个方法就是\textbf{Penrose 过程}（Penrose process）。想法很简单：从能层之外出发，你给自己装备一块大石头，跃向黑洞。如果我们把（你 $+$ 石头）这个系统的四维动量记为 $p^{(0)\mu}$，那么能量 $E^{(0)}=-K_{\mu}p^{(0)\mu}$ 无疑是正的，而且当你沿自己的测地线运动时它是守恒的。一旦你进入能层，就用尽全力、以非常特别的方式把石头猛掷出去。如果把你的动量记为 $p^{(1)\mu}$、石头的动量记为 $p^{(2)\mu}$，那么在你掷出石头的那一瞬间，正如狭义相对论中一样，我们有动量守恒：
\begin{equation*}
p^{(0)\mu} = p^{(1)\mu} + p^{(2)\mu}.
\tag{6.97}
\end{equation*}
与 Killing 矢量 $K_{\mu}$ 缩并，给出
\begin{equation*}
E^{(0)} = E^{(1)} + E^{(2)}.
\tag{6.98}
\end{equation*}
但是，如果我们设想你足够强壮（而且投掷得足够精准），你就可以按照式 (6.96) 来安排你的投掷，使得 $E^{(2)}<0$。此外，Penrose 还能够证明，你可以像图 6.9 那样安排初始轨迹和投掷，使得随后你沿一条测地线轨迹回到静止极限面之外、进入外部宇宙。由于你的能量沿途守恒，

%FIG{6.9}: {Penrose 过程（俯视图）。一个物体落向 Kerr 黑洞，并在能层之内（静止极限面以内、外侧事件视界之外）分裂为二。一块以负能量 $E^{(2)}$ 落入视界，另一块则以比原来落入物体更大的能量逃逸到无穷远。}

% ---- 书页 269 / p284 ----

最终我们将有
\begin{equation*}
E^{(1)} > E^{(0)}.
\tag{6.99}
\end{equation*}
这样，你脱身而出时携带的能量就比你进入时\emph{更多}。

世上没有免费的午餐；你获得的能量来自某处，而那个某处就是黑洞。事实上，Penrose 过程是通过减小转动黑洞的角动量来从中汲取能量的；要让这个戏法奏效，你必须逆着黑洞的转动方向掷出石头。为了更精确地看清这一点，回忆一下我们在本章前文中的论断：稳态时空中的任何事件视界，对于某个 Killing 矢量来说都是一个 Killing 视界。对 Kerr 而言，这个矢量是时间平移和转动 Killing 矢量的线性组合，
\begin{equation*}
\chi^{\mu} = K^{\mu} + \Omega_{\mathrm{H}}R^{\mu},
\tag{6.100}
\end{equation*}
其中 $\Omega_{\mathrm{H}}$ 恰恰就是式 (6.92) 中定义的视界角速度。利用 $K=\partial_t$ 和 $R=\partial_\phi$，不难验证 $\chi^{\mu}$ 在外侧事件视界处变为类光。具有动量 $p^{(2)\mu}$ 的粒子“在时间中向前”地穿越事件视界，这一陈述不过是
\begin{equation*}
p^{(2)\mu}\chi_{\mu} < 0.
\tag{6.101}
\end{equation*}
代入 $E$ 和 $L$ 的定义，我们看到这个条件等价于
\begin{equation*}
L^{(2)} < \frac{E^{(2)}}{\Omega_{\mathrm{H}}}.
\tag{6.102}
\end{equation*}
既然我们已经安排 $E^{(2)}$ 为负、$\Omega_{\mathrm{H}}$ 为正，这就意味着粒子必须带有负的角动量——它在逆着黑洞转动的方向运动。一旦你逃出能层、石头落入事件视界之内，黑洞的质量和角动量就是原来的数值再加上石头所带来的负的贡献：
\begin{gather*}
\delta M = E^{(2)}\\
\delta J = L^{(2)},
\tag{6.103}
\end{gather*}
其中 $J=Ma$ 是黑洞的角动量。于是式 (6.102) 就变成了对角动量能够减小多少的一个限制：
\begin{equation*}
\delta J < \frac{\delta M}{\Omega_{\mathrm{H}}}.
\tag{6.104}
\end{equation*}
如果我们恰好达到这个极限——随着我们掷入的石头越来越接近类光——我们就得到了“理想”过程，其中 $\delta J = \delta M/\Omega_{\mathrm{H}}$。

现在我们将利用这些想法来验证：尽管你可以用 Penrose 过程从黑洞中汲取能量（从而减小 $M$），你却不能

% ---- 书页 270 / p285 ----

违反面积定理：事件视界的面积是不减的。虽然质量减小了，但角动量也必须减小，而且两者减小的组合方式只允许面积增加。为了看清这一点，我们来计算外侧事件视界的面积，它位于
\begin{equation*}
r_{+} = GM + \sqrt{G^{2}M^{2}-a^{2}}.
\tag{6.105}
\end{equation*}
视界上的诱导度规 $\gamma_{ij}$（其中 $i$ 和 $j$ 取遍 $\{\theta,\phi\}$）可以直截了当地求得：在式 (6.70) 中令 $r=r_{+}$（于是 $\Delta=0$）、$\mathrm{d}t=0$ 和 $\mathrm{d}r=0$：
\begin{align*}
\gamma_{ij}\,\mathrm{d}x^{i}\mathrm{d}x^{j} &= ds^{2}(\mathrm{d}t=0,\,\mathrm{d}r=0,\,r=r_{+})\\
&= (r_{+}^{2}+a^{2}\cos^{2}\theta)\mathrm{d}\theta^{2} + \left[\frac{(r_{+}^{2}+a^{2})^{2}\sin^{2}\theta}{r_{+}^{2}+a^{2}\cos^{2}\theta}\right]\mathrm{d}\phi^{2}.
\tag{6.106}
\end{align*}
于是，视界面积就是诱导体积元的积分，
\begin{equation*}
A = \int\sqrt{|\gamma|}\,\mathrm{d}\theta\,\mathrm{d}\phi.
\tag{6.107}
\end{equation*}
行列式为
\begin{equation*}
|\gamma| = (r_{+}^{2}+a^{2})^{2}\sin^{2}\theta,
\tag{6.108}
\end{equation*}
所以视界面积就是
\begin{equation*}
A = 4\pi(r_{+}^{2}+a^{2}).
\tag{6.109}
\end{equation*}

为了证明面积不会减小，更方便的做法是改用黑洞的\textbf{不可约质量}（irreducible mass），其定义为
\begin{align*}
M_{\mathrm{irr}}^{2} &= \frac{A}{16\pi G^{2}}\\
&= \frac{1}{4G^{2}}(r_{+}^{2}+a^{2})\\
&= \frac{1}{2}\left(M^{2}+\sqrt{M^{4}-(Ma/G)^{2}}\right)\\
&= \frac{1}{2}\left(M^{2}+\sqrt{M^{4}-(J/G)^{2}}\right).
\tag{6.110}
\end{align*}
经过一点计算，我们可以通过求导得到 $M_{\mathrm{irr}}$ 如何受质量或角动量变化的影响：
\begin{equation*}
\delta M_{\mathrm{irr}} = \frac{a}{4GM_{\mathrm{irr}}\sqrt{G^{2}M^{2}-a^{2}}}\left(\Omega_{\mathrm{H}}^{-1}\delta M - \delta J\right).
\tag{6.111}
\end{equation*}

% ---- 书页 271 / p286 ----

于是，我们的极限 (6.104) 就变成
\begin{equation*}
\delta M_{\mathrm{irr}} > 0.
\tag{6.112}
\end{equation*}
不可约质量永远无法减小；名称由此而来。由此可知，在把黑洞的转动减缓到零之前，我们能够从黑洞汲取的最大能量为
\begin{equation*}
M - M_{\mathrm{irr}} = M - \frac{1}{\sqrt{2}}\left(M^{2}+\sqrt{M^{4}-(J/G)^{2}}\right)^{1/2}.
\tag{6.113}
\end{equation*}
这种彻底汲取的结果，是一个质量为 $M_{\mathrm{irr}}$ 的 Schwarzschild 黑洞。事实证明，我们所能做到的最好情形是从一个极端 Kerr 黑洞出发；这样我们可以取出它总能量的大约 29\%。

$M_{\mathrm{irr}}$ 的不可约性直接引出这样一个事实：面积 $A$ 永不减小。由式 (6.110) 和 (6.111) 我们有
\begin{equation*}
\delta A = 8\pi G\,\frac{a}{\Omega_{\mathrm{H}}\sqrt{G^{2}M^{2}-a^{2}}}\left(\delta M - \Omega_{\mathrm{H}}\delta J\right),
\tag{6.114}
\end{equation*}
它可以改写为
\begin{equation*}
\delta M = \frac{\kappa}{8\pi G}\delta A + \Omega_{\mathrm{H}}\delta J,
\tag{6.115}
\end{equation*}
其中我们引入了
\begin{equation*}
\kappa = \frac{\sqrt{G^{2}M^{2}-a^{2}}}{2GM\left(GM+\sqrt{G^{2}M^{2}-a^{2}}\right)}.
\tag{6.116}
\end{equation*}
量 $\kappa$ 当然就是 Kerr 解的表面引力，你可以把式 (6.100) 代入式 (6.9) 来验证这一点。

像 (6.115) 这样的方程，最先促使人们开始思考黑洞与热力学之间的对应。考虑热力学第一定律，
\begin{equation*}
\mathrm{d}E = T\,\mathrm{d}S - p\,\mathrm{d}V,
\tag{6.117}
\end{equation*}
其中 $T$ 是温度，$S$ 是熵，$p$ 是压强，$V$ 是体积，因此 $p\,\mathrm{d}V$ 项代表我们对系统所做的功。很自然地，我们可以把式 (6.115) 中的 $\Omega_{\mathrm{H}}\delta J$ 项看成我们向黑洞投掷石头时对它所做的功。这样，如果我们把热力学量——能量、熵和温度——与黑洞的质量、面积和表面引力等同起来，这个对应就开始成形：
\begin{gather*}
E \leftrightarrow M\\
S \leftrightarrow A/4G\\
T \leftrightarrow \kappa/2\pi.
\tag{6.118}
\end{gather*}

% ---- 书页 272 / p287 ----

（记住，我们使用的单位制中 $\hbar=c=k=1$。）在经典广义相对论的语境中，这个类比本质上是完美的：热力学的每一条定律都对应着一条黑洞力学定律。处于热平衡的系统将弛豫到一个稳态，对应于稳态黑洞。热力学第零定律说，在热平衡中温度在整个系统中处处相同；对黑洞的类似陈述是：稳态黑洞在整个视界上具有恒定的表面引力。至少在使事件视界成为 Killing 视界的那同样的合理假设之下，这一点将会成立。正如我们已经看到的，第一定律 (6.117) 等价于 (6.115)。第二定律——熵永不减小——不过就是视界面积永不减小这一陈述。最后，第三定律的通常表述是：在任何物理过程中都不可能达到 $T=0$，或者说，当温度趋于零时熵必须趋于零。对黑洞来说，这不太灵验；事实证明 $\kappa=0$ 对应于极端黑洞，而极端黑洞的面积不一定为零。不过，热力学第三定律本身其实也不真正灵验，因为确实存在违反它的普通物理系统；第三定律适用于某些情形，但并不是真正基本的定律。

在提出对应关系 (6.118) 时，我们有点儿作弊了；你会注意到，把 $T\mathrm{d}S$ 与 $\kappa\mathrm{d}A/8\pi G$ 等同起来之后，我们并不知道如何分别归一化 $S/A$ 或 $T/\kappa$，只知道它们的组合。不过，正如我们将在第 9 章讨论的，Hawking 证明了黑洞背景中的量子场使得黑洞以温度 $T=\kappa/2\pi$ 发出辐射。一旦知道了这一点，我们就可以把 $A/4G$ 解释为黑洞的真实熵。Bekenstein 提出了\textbf{广义第二定律}（generalized second law）：物质与黑洞合起来的熵永不减小：
\begin{equation*}
\delta\left(S + \frac{A}{4G}\right) \geq 0.
\tag{6.119}
\end{equation*}
广义第二定律实际上可以在多种不同的假设下得到证明。不过，我们通常喜欢把系统的熵与可到达的量子态数目的对数联系起来。因此，这个概念与无毛定理之间存在着某种张力：无毛定理告诉我们，具有固定电荷、质量和自旋的黑洞，其可能的状态极少（事实上只有一个）。这种行为似乎在向我们暗示，它是量子力学与引力之间相互作用之深刻性质的一个表征。

\section{习题}

\textbf{1.}\quad 证明：耦合的 Einstein–Maxwell 方程组可以同时由度规 (6.62) 和静电势 (6.67) 求解，只要 $H(\vec{x})$ 满足拉普拉斯方程
\begin{equation*}
\nabla^{2}H = 0.
\tag{6.120}
\end{equation*}

\textbf{2.}\quad 考虑 Kerr 黑洞赤道面内无质量粒子的轨道，仿射参量为 $\lambda$。

\textbf{(a)}\quad 证明
\begin{equation*}
\left(\frac{\mathrm{d}r}{\mathrm{d}\lambda}\right)^{2} = \frac{\Sigma^{2}}{\rho^{4}}\left(E - LW_{+}(r)\right)\left(E - LW_{-}(r)\right),
\tag{6.121}
\end{equation*}
其中 $\Sigma^{2} = (r^{2}+a^{2})^{2} - a^{2}\Delta(r)\sin^{2}\theta$，$E$ 和 $L$ 是守恒的能量和角动量，而 $W_{\pm}(r)$ 的表达式则要你自己求出。

\textbf{(b)}\quad 利用这个结果，并假设 $\Sigma^{2}$ 处处大于零，证明：赤道面内光子的轨道不可能在外侧事件视界 $r_{+}$ 之内存在转折点。这意味着，入射光线一旦穿越 $r_{+}$ 就无法逃逸，所以它确实是一个事件视界。

% ---- 书页 273 / p288 ----

\textbf{3.}\quad 在存在电磁场的情况下，电荷为 $e$、质量为 $m$ 的粒子服从
\begin{equation*}
\frac{\mathrm{d}^{2}x^{\mu}}{\mathrm{d}\tau^{2}} + \Gamma^{\mu}_{\rho\sigma}\frac{\mathrm{d}x^{\rho}}{\mathrm{d}\tau}\frac{\mathrm{d}x^{\sigma}}{\mathrm{d}\tau} = \frac{e}{m}F^{\mu}{}_{\nu}\frac{\mathrm{d}x^{\nu}}{\mathrm{d}\tau}.
\tag{6.122}
\end{equation*}
设想这样一个粒子在电荷为 $Q$、质量为 $M$ 的 Reissner–Nordström 黑洞的场中运动。

\textbf{(a)}\quad 证明能量
\begin{equation*}
E = m\left(1 - \frac{2GM}{r} + \frac{GQ^{2}}{r^{2}}\right)\frac{\mathrm{d}t}{\mathrm{d}\tau} + \frac{eQ}{r}
\tag{6.123}
\end{equation*}
是守恒的。

\textbf{(b)}\quad Penrose 型的过程对带电黑洞行得通吗？对于最大的物理过程，黑洞质量的改变量 $\delta M$ 是多少？

\textbf{4.}\quad 考虑静态坐标下的 de Sitter 空间：
\begin{equation*}
ds^{2} = -\left(1 - \frac{\Lambda}{3}r^{2}\right)\mathrm{d}t^{2} + \frac{\mathrm{d}r^{2}}{1 - \frac{\Lambda}{3}r^{2}} + r^{2}\mathrm{d}\Omega^{2}.
\end{equation*}
这个空间有一个 Killing 矢量 $\partial_t$，它在 $r=0$ 附近是类时的，而在某个 Killing 视界上是类光的。试定出 Killing 视界的径向位置 $r_{\mathrm{K}}$。该视界的表面引力 $\kappa$ 是多少？考虑通过替换 $t\rightarrow i\tau$ 得到的 de Sitter 空间的欧氏号差版本。证明：只要把 $\tau$ 取成周期的，就可以通过一个坐标变换使欧氏度规在视界处正则。

\textbf{5.}\quad 一位观者沿周长为 $2\pi R$ 的圆轨道，绕着电荷为 $Q$、质量为 $M$ 的 Reissner–Nordström 黑洞运行，他所看到的磁场是什么？

\textbf{6.}\quad 考虑一个 Kerr 黑洞，在其赤道面内有一个质量可忽略的吸积盘（accretion disk）。假设盘中的粒子沿测地线运动（也就是说，忽略任何压强支撑）。现在进一步假设盘中含有一些铁原子，它们正被某个辐射源激发。当铁原子退激发时，它们会发出辐射，在它们自己的静止系中测量，其频率为已知的 $\nu_{0}$。假设我们在远离黑洞的地方探测这种辐射（我们也处在赤道面内）。从盘的任一边缘以及从盘的中心发出的光子，其观测频率各是多少？考虑盘与黑洞同向转动和反向转动这两种情形。我们能否利用这些测量来确定黑洞的质量和角动量？
